\label{part:engineering}

\begin{introduction}
\item \textbf{Part IV}: Engineering Implementation - Morpho-Semantic Transformer
\item 从理论到工程的跳跃
\item MST 架构：形与质的计算机实现
\item 通往 AGI 的几何工程学
\end{introduction}

\textbf{(Part IV: Engineering - The Morpho-Semantic Transformer)}

这一部分是从理论转向工程实现的关键篇章。我们将从黑板公式转向 \textbf{MST (形质互变 Transformer)} 的可实现架构。

从双塔编码器的解耦，到融合注意力的耦合；从 T2I 生成，到具身智能交互，我们将逐步构建一套可执行的结构化引擎。这套引擎强调先搭建\textbf{拓扑骨架}，再填充\textbf{纤维内容}。

目标不仅是让机器生成“像苹果”的图像，更是让它在内部表示中捕捉“圆”与“红”如何几何耦合。换言之，我们希望它生成的不只是外观相似样本，而是\textbf{构造上更一致}的对象表示。


%--chp---
